\documentclass[11pt,a4paper]{article}
\usepackage[margin=19mm,headheight=15pt]{geometry}
\usepackage{fontspec}
\setmainfont{Noto Serif}
\usepackage{amsmath,amssymb,mathtools}
\usepackage{xcolor}
\usepackage{fancyhdr}
\usepackage{enumitem}
\usepackage{microtype}
\setlength{\parindent}{0pt}
\setlength{\parskip}{6pt}
\definecolor{navy}{HTML}{183153}
\definecolor{lightgray}{HTML}{F3F5F7}
\pagestyle{fancy}\fancyhf{}
\lhead{\small ME 5309 -- Mid-Semester Examination}
\rhead{\small ME 5309}
\cfoot{\thepage}
\begin{document}
\begin{center}
{\Large\bfseries ME 5309 -- Mid-Semester Examination}\\[3pt]
{\large Answer Booklet}\\[8pt]
\end{center}
\textbf{Instructions:} Attempt any three questions. Each question carries 10 marks. Show working where appropriate.
\section*{Question 1 \hfill [10 marks]}
\subsection*{Q1(a) \hfill [4 marks]}

The pixel pitch is 0.005 mm, so f\_x = 16/0.005 = 3200 pixels.
Sensor width = 6.4 mm.
Belt field of view = 400x6.4/16 = 160 mm; spatial resolusion = 160/1280 = 0.125 mm/pixel.

\subsection*{Q1(b) \hfill [3 marks]}

ADC step = 5/2\textasciicircum{}12 = 1.2207 mV.
Since sensitivity is 10 mV/N, ΔF = 1.2207/10 = 0.122 N.
I report this as 0.122 mN.

\subsection*{Q1(c) \hfill [3 marks]}

For zero-order hold, F = e\textasciicircum{}(AT) and G = integral\_0\textasciicircum{}T e\textasciicircum{}(Aτ)B dτ.
Here:\[\begin{aligned}A = -2, B = 1, T = 0.1 s\end{aligned}\]
Thus:\[\begin{aligned}F = e^(-0.2) = 0.818731\\ G = (1-e^(-0.2))/2 = 0.0906346\end{aligned}\]
Therefore:\[\begin{aligned}x_{k+1} = 0.818731 x_k + 0.0906346 u_k\end{aligned}\]\par\smallskip With $AT=(-2)(0.1)=-0.2$, $F=e^{-0.2}=0.8187308$. For $G$, $(1-F)/2=(1-0.8187308)/2=0.1812692/2=0.0906346$. Thus the input contribution is positive.

\section*{Question 2 \hfill [10 marks]}
\subsection*{Q2(a) \hfill [2 marks]}

Choose x₁ = θ and x₂ = θ̇.
Then ẋ₁ = x₂ and ẋ₂ = [τ - b x₂ - mgℓ sin(x₁)]/J.
At equilibrium:\[\begin{aligned}x_2 = 0 and τ_0 = mg\ell sin \theta_e\end{aligned}\]
Since |sin θ\_e| less than or equal to 1, an equilibrium exists when |τ₀| less than or equal to mgℓ = 2.4525 N m.\par\smallskip At rest $\dot\theta=0$ and $\ddot\theta=0$, so the damping term vanishes. Rearranging gives $\tau_0=mg\ell\sin\theta_e$. Here $mg\ell=1\times9.81\times0.25=2.4525$ N\,m. Since $-1\leq\sin\theta_e\leq1$, the required torque must lie between $-2.4525$ and $+2.4525$ N\,m.

\subsection*{Q2(b) \hfill [2 marks]}

With x = [θ, θ̇]ᵀ, A = [[0,1], [-(mgℓ/J)cos θ\_e, -b/J]], B = [[0],[1/J]].
At θ\_e = 60°, A = [[0,1], [-61.3125, -2.5]] and B = [[0],[50]].\par\smallskip At the equilibrium, the constant torque balances gravity. For a small perturbation, use $\sin(\theta_e+\delta\theta)\approx\sin\theta_e+\cos\theta_e\,\delta\theta$. The constant terms cancel against $\tau_0$; the remaining linear term gives the local state matrix.

\subsection*{Q2(c) \hfill [4 marks]}

At rest, θ̇ = θ̈ = 0, so the torque balance is Kp(θd-θss) = mgℓ sin θss.
With Kp = 50, θd = π/2 and mgℓ = 2.4525, solve 50(π/2-θss) = 2.4525 sin θss.
This gives θss approximately 1.52181 rad = 87.193°, so the error is 90°-87.193° = 2.807°; it is not zero.\par\smallskip Rearranging the steady-state equation gives $50(\pi/2-\theta_{ss})=2.4525\sin\theta_{ss}$. I used radians in the equation. Solving numerically gives $\theta_{ss}\approx1.522\,$rad $\approx87.19^\circ$, so the gap from $90^\circ$ is about $2.81^\circ$.

\subsection*{Q2(d) \hfill [2 marks]}

\vspace{2mm}

\section*{Question 3 \hfill [10 marks]}
\subsection*{Q3(a) \hfill [2 marks]}

Finite differencing amplifies encoder noise, making the velocity estimate noisy.

\subsection*{Q3(b) \hfill [2 marks]}

\[\begin{aligned}y[k]=0.8818y[k-1]+0.0591x_k+0.0591x[k-1], DC gain 1\end{aligned}\]

\subsection*{Q3(c) \hfill [4 marks]}

With L=[ℓ₁,ℓ₂]ᵀ and C=[1,0], A-LC=[[-ℓ₁,1], [-ℓ₂,-2]].
Its characteristic polynomial is det(sI-(A-LC)) = (s+ℓ₁)(s+2)+ℓ₂ = s²+(ℓ₁+2)s+(2ℓ₁+ℓ₂).
Desired poles at -10 give (s+10)²=s²+20s+100.
Hence ℓ₁+2=20 gives ℓ₁=18, and 2ℓ₁+ℓ₂=100 gives ℓ₂=64.\par\smallskip Expanding $(s+\ell_1)(s+2)+\ell_2$ gives $s^2+(\ell_1+2)s+(2\ell_1+\ell_2)$. The desired polynomial is $(s+10)^2=s^2+20s+100$. Matching the $s$ coefficient gives $\ell_1=18$; matching the constant term then gives $36+\ell_2=100$, hence $\ell_2=64$.

\subsection*{Q3(d) \hfill [2 marks]}

An observer can estimate velocity from position measurements.
A filter can be safer because it smooths noise.

\section*{Question 4 \hfill [10 marks]}
\subsection*{Q4(a) \hfill [2 marks]}

Closed-loop characteristic polynomial is s²+(a+bk₂)s+(c+bk₁)=s²+4s+8, giving poles -2±2j.
At DC, θ/r = bNff/(c+bk₁), so unity gain requires Nff=(c+bk₁)/b=(6+2)/1=8.\par\smallskip Comparing coefficients with $s^2+4s+8$ gives $a+bk_2=4$ and $c+bk_1=8$. The feedforward gain is chosen from the nominal plant so the reference scaling is unity; here this gives $N_{ff}=8$. The closed-loop roots of $s^2+4s+8$ are $-2\pm2j$.

\subsection*{Q4(b) \hfill [2 marks]}

\[\begin{aligned}\theta_{ss} = 0.75\times8/(6+2)=0.75,\\ \text{the percentage tracking error is 25\%}\end{aligned}\]

\subsection*{Q4(c) \hfill [4 marks]}

The characteristic polynomial is s³+(a+bk₂)s²+(c+bk₁)s+bK\_I.
I use the target s³+6s²+12s+16, giving k₂=4, k₁=6, K\_I=16.

\subsection*{Q4(d) \hfill [2 marks]}

\vspace{2mm}

\end{document}